\section{Discussion and Future Work}
\label{sec:discussion}

% Risk table (Section VII-A) defined here so the full-width float lands before the references.
\begin{table*}[t]
  \centering
  \caption{Updated risk and ethics reflection: anticipated (A) and newly emerged (N) issues and how each was handled.}
  \label{tab:risks}
  \footnotesize
  \begin{tabular}{@{}p{2.9cm} p{6.7cm} p{7.2cm}@{}}
    \toprule
    Issue & What emerged & How it was managed \\
    \midrule
    Class imbalance (A) &
      Worse than expected: 6,382:1 between \emph{tree} and \emph{sign};
      safety classes 1.1\% of training pixels. &
      Four imbalance-aware losses vs.\ CE; train-only median-frequency weights,
      clipped at 50; per-class IoU reported. \\
    Split shift (N) &
      \emph{Bridge} only in test; validation lacks four classes; test is
      hazard-heavy (safety classes 10.3\% vs.\ 1.1\%), driven by \emph{creek};
      val closer to train than to test. &
      Official split kept for comparability; \emph{bridge} excluded from headline
      mIoU but reported; val and test reported separately; test per-class IoU;
      test guarded by a flag and used once. \\
    Leakage and labels (A) &
      Labeled frames are near-duplicates (every fifth video frame); coarse
      polygons, wires labeled \emph{tree}, sticks labeled \emph{log},
      near-identical concrete/rock-bed colors. &
      Sequence-level split verified (20-frame discrepancy with~\cite{wigness2019rugd}
      disclosed); boundary-band analysis (Section~\ref{sec:error}); palette
      caveat in captions. \\
    Compute (A) &
      Apple M1 Pro too slow at $512^2$ and unstable in validation; Colab T4
      sessions capped at about 12\,h. &
      $384^2$/25 epochs for \emph{all} runs; mixed precision on T4; per-epoch
      checkpoints on Drive with automatic resume (Section~\ref{sec:compute}).
      Outcome: quota ran out before 3 of 4 seed reruns, and four checkpoints
      were lost to a storage mishap; disclosed, and the selected configuration's
      replicate used for test. \\
    Domain shift (A) &
      Different vehicle, viewpoint, and label mix in RELLIS-3D; \emph{mud} has
      no RUGD counterpart. &
      Group map fixed before evaluation; mud ignored in the headline score with
      a sensitivity check; zero-shot only. Outcome: group mIoU fell from 60.0\% to
      43.4\%; ground and water transfer worst (Section~\ref{sec:res_rellis}). \\
    Safety, dual use, licensing &
      Rare hazards are the most likely misses; Army-funded data;
      RELLIS-3D is CC BY-NC-SA 3.0. &
      Rare-class IoU reported; models presented as research artifacts;
      non-commercial use; no redistribution of images. \\
    \bottomrule
  \end{tabular}
\end{table*}

\subsection{What Worked}

Adding a Dice term to the loss was the most useful single change: about 3
points of GT-present validation mIoU over CE (Section~\ref{sec:res_pretrain}),
concentrated on water, log, container, and concrete, with frequent classes
and pixel accuracy unchanged. Because Dice normalizes per class over the
batch, a rare class's gradient no longer scales with its pixel count, which
is the imbalance a pixel-averaged loss suffers from. ImageNet pretraining was
the largest single factor (about 6 points) and mattered most for the same
rare objects. The DeepLabV3+ decoder helped moderately (about 2 points),
mainly on objects whose extent needs multi-scale context, such as vehicles
and picnic tables.

\subsection{What Did Not Work}

Median-frequency weighted CE hurt across the board ($-9.6$ validation mIoU,
$-5.0$ pixel accuracy): weights up to 50 made the model over-predict rare
classes, so their false positives grew faster than their true positives.
Deepening the encoder from ResNet-18 to ResNet-50 bought less than seed
noise. \emph{Dirt}, \emph{sign}, and \emph{person} stayed near zero under
every loss: dirt is 0.13\% of training pixels and visually close to mulch and
sand, and signs and people are tiny at $384^2$. Best-epoch selection on a
noisy, incomplete validation split was itself a weakness. Under our
pre-registered rule, Dice alone looked 2 points worse than CE+Dice, when on
the classes validation actually contains the two are tied. The validation-to-test gap (7.5 points for the final model) was
concentrated in one scene: creek, whose clear stream the model reads as mulch,
caused 80\% of test errors, and without it test mIoU (49.3\%) matches
validation. A validation split with no flowing water could not warn us. Some limitations
are structural rather than methodological. Validation lacks four classes, so
model selection cannot reward improvements on \emph{sand}, \emph{rock-bed},
\emph{bicycle}, or \emph{bridge}, and its class mix is closer to training
than to test, so it is an optimistic proxy. The test split's rare hazards come
almost entirely from one sequence, which makes their IoU a high-variance
estimate. The reduced $384^2$, 25-epoch schedule
(Section~\ref{sec:compute}) limits thin-structure accuracy and may understate
what each configuration can reach. With two seeds for one
configuration (the other planned reruns did not fit the free GPU quota),
we can bound run-to-run noise only roughly, and the test results for the
selected configuration come from its seed-1 replicate.

\subsection{Future Work}

(i)~Train at native resolution or with sliding-window inference to recover
thin obstacles. (ii)~Re-split or use sequence-level cross-validation so that
validation covers every trainable class and matches the test class mix.
(iii)~Exploit temporal consistency across RUGD's video frames. (iv)~Compare
transformer backbones and group-aware objectives such as
GA-Nav~\cite{guan2022ganav}. (v)~Calibrate uncertainty so a planner can treat
low-confidence water or obstacle pixels conservatively. (vi)~Fine-tune on a
small RELLIS-3D subset, drawn by sequence rather than by frame, to measure how
much target data closes the cross-dataset gap, which was largest for ground
and water (Section~\ref{sec:res_rellis}). (vii)~Add water-specific cues, such
as specular-reflection or temporal-flicker features, or a small labeled set
of streams, since appearance alone failed on clear, shallow water.
